\section{Análisis de la Suma Ponderada (Pregunta 2)}

\textit{¿Existe una diferencia significativa entre la solución encontrada
calculando los pesos a priori con el GA frente a seleccionarla a posteriori
del frente de Pareto generado por NSGA-II? Explica por qué coinciden o
difieren.}

\subsection{Comparación cuantitativa}

\begin{table}[h]
  \centering
  \caption{Suma ponderada: A Priori (GA) vs.\ A Posteriori (NSGA-II).}
  \label{tab:ponderada}
  \begin{tabular}{lSSS}
    \toprule
    & {\textbf{$f_1$ (Costo)}} & {\textbf{$f_2$ (Tiempo)}} & {\textbf{Suma escalada}} \\
    \midrule
    A Priori (GA)        & \AprioriPondCosto & \AprioriPondTiempo & {---} \\
    A Posteriori (NSGA-II)& \PostPondCosto   & \PostPondTiempo    & {---} \\
    \midrule
    Diferencia absoluta  & \DiffPondCosto    & \DiffPondTiempo    & {---} \\
    \bottomrule
  \end{tabular}
\end{table}

Las dos soluciones son \textbf{muy cercanas pero no idénticas}: difieren en
apenas $\num{\DiffPondCosto}$ unidades de costo y $\num{\DiffPondTiempo}$
unidades de tiempo. En la Figura~\ref{fig:scatter} ambos marcadores (estrella
oro y equis roja) caen prácticamente superpuestos sobre el ``codo'' del frente.

\subsection{Por qué coinciden}

La razón es geométrica. El método de suma ponderada equivale a buscar la
solución que toca primero una recta de pendiente $-w_1/w_2$ que barre el
espacio de objetivos. Cuando el frente de Pareto es \textbf{convexo} ---como
ocurre aquí--- esa recta toca el frente en un único punto bien definido, y
\emph{cualquier} método que minimice la misma combinación lineal converge
esencialmente al mismo lugar~\cite{marler2004survey,miettinen1999nonlinear}.
Por construcción, el problema $f_1 = 4(x_1^2+x_2^2)$, $f_2 = (x_1-5)^2+(x_2-5)^2$
genera un frente convexo, lo que explica la coincidencia.

De hecho, la solución A Priori converge a $\mathbf{x} \approx
(\num{\AprioriPondVarUno}, \num{\AprioriPondVarDos})$, que coincide con el
óptimo analítico exacto $(1.5, 1.5)$ obtenible por simetría e igualación de
derivadas. El GA recuperó el óptimo teórico sin error apreciable.

\subsection{Por qué difieren levemente}

La pequeña discrepancia se debe a que \textbf{ambos enfoques no minimizan
exactamente la misma función escalar}:

\begin{itemize}[noitemsep]
  \item El A Priori usa un \textbf{escalamiento manual fijo} ($f_1/200$,
        $f_2/50$), denominadores elegidos a ojo según los rangos esperados.
  \item El A Posteriori \textbf{normaliza por los extremos reales} del frente
        muestreado, que dependen de las \MetNSol{} soluciones encontradas.
\end{itemize}

Como los factores de escala difieren, la dirección de preferencia efectiva se
inclina ligeramente distinto, desplazando el punto seleccionado unas pocas
unidades. Además, el A Posteriori sólo puede elegir entre los puntos
\emph{discretos} del frente muestreado, no sobre el continuo. En resumen: no
hay diferencia \emph{significativa} ---el sesgo del escalamiento manual es
marginal cuando el frente es convexo--- pero sí una diferencia \emph{medible}
y explicable.
